\documentclass[12pt,a4paper]{article}
\usepackage[utf8]{inputenc}
\usepackage[T5]{fontenc}
\usepackage[vietnamese]{babel}
\usepackage[margin=1.3cm,top=1.0cm,bottom=1.05cm,includehead,headheight=18pt,headsep=8pt,footskip=15pt]{geometry}
\usepackage{amsmath,amssymb,array,tabularx,enumitem,fancyhdr,tikz,xcolor,needspace}
\usetikzlibrary{calc,angles,quotes,arrows.meta,babel,patterns}
\definecolor{ink}{RGB}{30,55,70}

\newcommand{\headerQuote}{%
  \ifcase\value{page}\relax
    \or Math gives us hope that every problem has a solution% Trang 1
    \or Math + You = Champion% Trang 2
    \or Where Frustration Becomes Inspiration% Trang 3
    \or Turning Confusion into Clarity% Trang 4
    \or Don't be an X, be a Y of mathematics% Trang 5
    \or Don't get mad. Learn to add% Trang 6
    \or The only way to learn mathematics is to do mathematics% Trang 7
    \or Small steps every day lead to big results% Trang 8
    \or Every expert was once a beginner% Trang 9
    \or Success is the sum of small efforts repeated daily% Trang 10
    \or Mathematics is the music of reason% Trang 11
    \or Believe you can and you're halfway there% Trang 12
    \else Math gives us hope that every problem has a solution%
  \fi
}
\pagestyle{fancy}\fancyhf{}
\lhead{\small\textit{\headerQuote}}
\rhead{\small\textit{Người biên soạn: Dương Gia Bảo}}
\rfoot{\small Trang \thepage}
\renewcommand{\headrulewidth}{0.4pt}
\setlength{\headheight}{18pt}
\setlength{\parindent}{0pt}
\setlength{\parskip}{0pt}
\setlength{\tabcolsep}{1.5pt}
\renewcommand{\arraystretch}{1.0}
\newcommand{\parttitle}[1]{\par\vspace{1.5pt}\needspace{2\baselineskip}{\large\bfseries\color{ink}#1}\par\nopagebreak\vspace{1pt}}
\newcommand{\four}[4]{\par\noindent\begin{tabularx}{\linewidth}{@{}*{4}{>{\raggedright\arraybackslash}X}@{}}\textbf{A.} #1&\textbf{B.} #2&\textbf{C.} #3&\textbf{D.} #4\end{tabularx}}
\newcommand{\two}[4]{\par\noindent\begin{tabularx}{\linewidth}{@{}*{2}{>{\raggedright\arraybackslash}X}@{}}\textbf{A.} #1&\textbf{B.} #2\\[1pt]\textbf{C.} #3&\textbf{D.} #4\end{tabularx}}
\newcommand{\answerbox}{\hfill\textit{Đáp số:}\ \fbox{\rule{0pt}{9pt}\hspace{2.2cm}}}
\setlist[enumerate,1]{label=\textbf{Câu \arabic*.},leftmargin=1.1cm,labelsep=.1cm,itemsep=0.2pt,topsep=0.2pt,parsep=0pt}
\setlist[enumerate,2]{label=\alph*),leftmargin=.55cm,itemsep=0.2pt,topsep=0.2pt,parsep=0pt}

\newcommand{\soDe}{27}

\begin{document}
\noindent
\begin{minipage}[c]{0.38\textwidth}
    \centering
    \textbf{LÀM TOÁN ĐỂ HẠNH PHÚC}\\[-2pt]
    \rule{3.2cm}{0.6pt}\\[2pt]
    \textbf{ĐỀ CHÍNH THỨC}
\end{minipage}\hfill
\begin{minipage}[c]{0.35\textwidth}
    \centering
    \textbf{ĐỀ ÔN THI GIỮA KÌ I}\\
    \textbf{NĂM HỌC 2026 -- 2027}\\[1pt]
    \textbf{ĐỀ SỐ \soDe}\\
    \textbf{Môn: TOÁN -- Khối lớp 10}\\
    \textit{\small Thời gian làm bài: 90 phút\\(không kể thời gian phát đề)}
\end{minipage}

\vspace{0.04cm}
\noindent
\begin{minipage}[b]{0.78\textwidth}
    \textbf{Họ và tên:}\dotfill\textbf{SBD:}\makebox[3.5cm]{\dotfill}
\end{minipage}\hfill
\begin{minipage}[b]{0.20\textwidth}
    \hfill
    \setlength{\fboxrule}{0.8pt}%
    \fbox{\begin{minipage}{2.3cm}
        \centering
        \textbf{\small Mã đề thi}\\
        \textbf{00\soDe}
    \end{minipage}}
\end{minipage}
\vspace{0.04cm}

% =========================================================================
% TRANG 1: PHẦN I (12 CÂU TRẮC NGHIỆM NHIỀU LỰA CHỌN)
% Quy chuẩn: 12 câu nằm trọn vẹn trong Trang 1, kết thúc bằng \newpage.
% Cân bằng đáp án: 3A - 3B - 3C - 3D.
% =========================================================================
\parttitle{PHẦN I. Trắc nghiệm nhiều lựa chọn}
\textit{Thí sinh trả lời từ câu 1 đến câu 12. Mỗi câu hỏi chỉ chọn một phương án đúng duy nhất.}
\begin{enumerate}
\item Trong các câu sau, câu nào là một mệnh đề toán học?
\two{$17$ là số nguyên tố.}{Thời tiết hôm nay đẹp quá!}{Bạn có thích học môn Toán không?}{Hãy làm bài tập về nhà trước khi đến lớp!}

\item Phủ định của mệnh đề $P\colon \text{`}\forall x \in \mathbb{R},\, x^2 - 3x + 2 > 0\text{'}$ là
\two{$\overline{P}\colon \text{`}\forall x \in \mathbb{R},\, x^2 - 3x + 2 \le 0\text{'}$.}{$\overline{P}\colon \text{`}\exists x \in \mathbb{R},\, x^2 - 3x + 2 \le 0\text{'}$.}{$\overline{P}\colon \text{`}\exists x \in \mathbb{R},\, x^2 - 3x + 2 < 0\text{'}$.}{$\overline{P}\colon \text{`}\forall x \in \mathbb{R},\, x^2 - 3x + 2 < 0\text{'}$.}

\item Cho tập hợp $A = \{x \in \mathbb{N} \mid (2x - 1)(x^2 - 4) = 0\}$. Viết tập hợp $A$ bằng cách liệt kê các phần tử.
\four{$A = \left\{-2;\, \frac{1}{2};\, 2\right\}$.}{$A = \{-2;\, 2\}$.}{$A = \{2\}$.}{$A = \varnothing$.}

\item Cho tập hợp $X = \{a;\, b;\, c\}$. Số tập hợp con của tập hợp $X$ là
\four{$3$.}{$6$.}{$7$.}{$8$.}

\item Bất phương trình nào sau đây là bất phương trình bậc nhất hai ẩn?
\four{$2x - 3y + 5 > 0$.}{$x^2 + 2y \le 4$.}{$2x - 3y^2 + 1 \ge 0$.}{$xy + x - 2 < 0$.}

\item Điểm nào sau đây thuộc miền nghiệm của bất phương trình $2x - y + 3 \ge 0$?
\four{$(1;\, 6)$.}{$(2;\, 1)$.}{$(0;\, 4)$.}{$(-3;\, 1)$.}

\item Cặp số $(x;\, y)$ nào sau đây là một nghiệm của hệ bất phương trình $\begin{cases} x + y - 2 \le 0 \\ 2x - 3y + 6 > 0 \end{cases}$?
\four{$(2;\, 2)$.}{$(1;\, 3)$.}{$(0;\, 1)$.}{$(-2;\, 5)$.}

\item Đường thẳng nào dưới đây là bờ của nửa mặt phẳng miền nghiệm của bất phương trình $3x - 2y + 6 \ge 0$?
\two{$d\colon 2x - 3y + 6 = 0$.}{$d\colon 3x + 2y - 6 = 0$.}{$d\colon 3x - 2y - 6 = 0$.}{$d\colon 3x - 2y + 6 = 0$.}

\item Giá trị của $\cos 120^\circ$ bằng
\four{$-\dfrac{1}{2}$.}{$\dfrac{1}{2}$.}{$-\dfrac{\sqrt{3}}{2}$.}{$\dfrac{\sqrt{3}}{2}$.}

\item Cho tam giác $ABC$ có $AB = 5$, $AC = 8$ và $\widehat{A} = 60^\circ$. Độ dài cạnh $BC$ bằng
\four{$\sqrt{129}$.}{$7$.}{$\sqrt{41}$.}{$8$.}

\item Cho tam giác $ABC$ có $a = 6$, $b = 8$, $c = 10$. Bán kính $R$ của đường tròn ngoại tiếp tam giác $ABC$ bằng
\four{$4$.}{$2$.}{$5$.}{$10$.}

\item Cho tam giác $ABC$ có độ dài ba cạnh là $a, b, c$. Công thức tính độ dài đường trung tuyến $m_a$ kẻ từ đỉnh $A$ là
\two{$m_a^2 = \dfrac{b^2 + c^2 - a^2}{2}$.}{$m_a^2 = \dfrac{2(b^2 + c^2) + a^2}{4}$.}{$m_a^2 = \dfrac{b^2 + c^2}{2} - \dfrac{a^2}{2}$.}{$m_a^2 = \dfrac{2(b^2 + c^2) - a^2}{4}$.}
\end{enumerate}
\newpage

% =========================================================================
% TRANG 2: PHẦN II (2 CÂU ĐÚNG/SAI) + PHẦN III (CÂU 1 & CÂU 2 CÓ HÌNH TIKZ)
% Quy chuẩn: Nằm trọn vẹn trong Trang 2, kết thúc bằng \newpage.
% =========================================================================
\parttitle{PHẦN II. Trắc nghiệm đúng/sai}
\textit{Thí sinh trả lời từ câu 1 đến câu 2. Trong mỗi ý a), b), c), d) ở mỗi câu, thí sinh chọn đúng hoặc sai.}
\begin{enumerate}
\item Cho hai tập hợp $A = [-2;\, 4]$ và $B = (1;\, +\infty)$.
\begin{enumerate}
    \item $0 \in A$.
    \item Tập hợp $B$ là nửa khoảng $[1;\, +\infty)$.
    \item $A \cap B = (1;\, 4]$.
    \item $A \setminus B = [-2;\, 1]$.
\end{enumerate}

\item Cho tam giác $ABC$ có độ dài các cạnh $a = 7$, $b = 5$ và $c = 6$.
\begin{enumerate}
    \item Nửa chu vi của tam giác $ABC$ là $p = 9$.
    \item Góc $A$ của tam giác $ABC$ là một góc nhọn.
    \item $\cos A = \dfrac{1}{5}$.
    \item Diện tích của tam giác $ABC$ bằng $12\sqrt{6}$.
\end{enumerate}
\end{enumerate}

\parttitle{PHẦN III. Trắc nghiệm trả lời ngắn}
\textit{Thí sinh trả lời từ câu 1 đến câu 4. Điền kết quả vào ô tương ứng.}
\begin{enumerate}
\item Cho đường thẳng $d\colon x + 2y - 4 = 0$ chia mặt phẳng tọa độ $Oxy$ thành hai nửa mặt phẳng (hình vẽ). Miền không bị gạch chéo (kể cả bờ $d$) là miền nghiệm của bất phương trình bậc nhất hai ẩn có dạng $x + by \le c$. Tính giá trị của biểu thức $T = b + c$.
\begin{center}
\vspace{-3pt}
\begin{tikzpicture}[scale=0.72, >=Stealth]
    \draw[line width=1.1pt, ->] (-1.3,0) -- (5.2,0) node[right, font=\small\bfseries] {$x$};
    \draw[line width=1.1pt, ->] (0,-0.8) -- (0,3.2) node[above, font=\small\bfseries] {$y$};
    \node[below left=2pt, font=\small\bfseries] at (0,0) {$O$};
    
    % Miền bị gạch: nửa mặt phẳng x + 2y > 4
    \fill[pattern=north east lines, pattern color=gray!60] (-1.0, 2.5) -- (5.0, 2.5) -- (5.0, -0.5) -- cycle;
    \draw[line width=1.3pt, ink] (-1.0, 2.5) -- (5.0, -0.5);
    \node[above right=2pt, font=\small\bfseries] at (-1.0, 2.5) {$d$};
    
    \filldraw[black] (4,0) circle (2.4pt) node[below=2pt, font=\small\bfseries] {$4$};
    \filldraw[black] (0,2) circle (2.4pt) node[left=2pt, font=\small\bfseries] {$2$};
\end{tikzpicture}
\vspace{-6pt}
\end{center}
\answerbox

\item Để đo khoảng cách giữa hai vị trí $A$ và $B$ ở hai bên bờ một khúc sông mà không thể qua sông trực tiếp, một người quan sát chọn vị trí $C$ cùng bên bờ với điểm $A$ sao cho $AC = 60\text{ m}$. Người quan sát đo được các góc $\widehat{ACB} = 75^\circ$ và $\widehat{CAB} = 60^\circ$ (hình vẽ). Tính khoảng cách giữa hai điểm $A$ và $B$ theo đơn vị mét (làm tròn kết quả đến hàng đơn vị).
\begin{center}
\vspace{-3pt}
\begin{tikzpicture}[scale=0.82, >=Stealth]
    % Tọa độ các điểm: AC = 3.6 cm (tương ứng 60m), góc A = 60 độ, góc C = 75 độ
    \coordinate (A) at (0,0);
    \coordinate (C) at (3.6,0);
    \coordinate (B) at (2.459, 4.259);
    
    % Dòng sông minh họa bằng dải màu xanh mờ
    \fill[teal!12] (-0.8, 1.2) -- (4.4, 1.2) -- (4.4, 3.4) -- (-0.8, 3.4) -- cycle;
    \draw[teal!60, line width=0.8pt, dashed] (-0.8, 1.2) -- (4.4, 1.2);
    \draw[teal!60, line width=0.8pt, dashed] (-0.8, 3.4) -- (4.4, 3.4);
    \node[teal!80!black, font=\footnotesize\itshape] at (0.3, 2.3) {Dòng sông};
    
    % Tam giác ABC
    \draw[line width=1.3pt, ink] (A) -- (C) node[midway, below=3pt, font=\small\bfseries] {$60\text{ m}$};
    \draw[line width=1.3pt, ink] (C) -- (B);
    \draw[line width=1.3pt, dashed, red!85!black] (A) -- (B) node[midway, above left=1pt, font=\small\bfseries] {$AB = ?$};
    
    % Các điểm mút
    \filldraw[black] (A) circle (2.4pt) node[below left=2pt, font=\small\bfseries] {$A$};
    \filldraw[black] (C) circle (2.4pt) node[below right=2pt, font=\small\bfseries] {$C$};
    \filldraw[black] (B) circle (2.4pt) node[above=2pt, font=\small\bfseries] {$B$};
    
    % Cung góc A (60 độ)
    \draw[line width=0.9pt] (0.6,0) arc (0:60:0.6);
    \node[font=\small\bfseries] at (0.9, 0.45) {$60^\circ$};
    
    % Cung góc C (75 độ)
    \draw[line width=0.9pt] (3.0,0) arc (180:105:0.6);
    \node[font=\small\bfseries] at (2.7, 0.5) {$75^\circ$};
\end{tikzpicture}
\vspace{-6pt}
\end{center}
\answerbox
\newpage

% =========================================================================
% TRANG 3: PHẦN III (CÂU 3 & CÂU 4) + PHẦN IV (4 CÂU TỰ LUẬN)
% Quy chuẩn: Nằm trọn vẹn trong Trang 3, kết thúc bằng \newpage.
% =========================================================================
\item Một lớp học có $45$ học sinh. Trong một buổi khảo sát về thể thao, có $28$ học sinh thích môn Bóng đá, $23$ học sinh thích môn Bóng rổ và có $6$ học sinh không thích cả hai môn này. Hỏi có bao nhiêu học sinh trong lớp thích cả hai môn Bóng đá và Bóng rổ?
\answerbox

\item Cho tam giác $ABC$ có độ dài các cạnh $AB = 3$, $AC = 5$ và góc $\widehat{A} = 120^\circ$. Tính độ dài cạnh $BC$.
\answerbox
\end{enumerate}

\parttitle{PHẦN IV. Tự luận}
\textit{Thí sinh trình bày chi tiết lời giải từ câu 1 đến câu 4.}
\begin{enumerate}
\item \textit{(0,75 điểm)} Cho góc $\alpha$ thỏa mãn $\cos \alpha = -\dfrac{3}{5}$ với $90^\circ < \alpha < 180^\circ$. Tính giá trị của biểu thức
\[
P = \dfrac{2\sin \alpha - 3\cos \alpha}{\sin \alpha + \cos \alpha}.
\]

\item \textit{(0,75 điểm)} Cho hai tập hợp $A = [-3;\, 2)$ và $B = (0;\, 5]$.
\begin{enumerate}
    \item Xác định các tập hợp $A \cup B$, $A \cap B$ và $A \setminus B$.
    \item Biểu diễn tập hợp $A \cap B$ trên trục số.
\end{enumerate}

\item \textit{(0,75 điểm)} Cho hai tập hợp khác rỗng $A = [m - 1;\, m + 3]$ và $B = (-2;\, 4)$. Tìm tất cả các giá trị thực của tham số $m$ để tập hợp $A$ là tập con của tập hợp $B$ ($A \subset B$).

\item \textit{(0,75 điểm)} Một xưởng cơ khí dự định sản xuất hai loại sản phẩm $A$ và $B$. Để sản xuất một đơn vị sản phẩm loại $A$ cần $2\text{ kg}$ thép và $1$ giờ làm việc của máy. Để sản xuất một đơn vị sản phẩm loại $B$ cần $1\text{ kg}$ thép và $2$ giờ làm việc của máy. Xưởng hiện có sẵn $10\text{ kg}$ thép và máy có thể hoạt động tối đa $8$ giờ mỗi ngày. Mỗi đơn vị sản phẩm loại $A$ đem lại lợi nhuận $400$ nghìn đồng, mỗi đơn vị sản phẩm loại $B$ đem lại lợi nhuận $300$ nghìn đồng. Giả sử toàn bộ sản phẩm xưởng sản xuất trong ngày đều được tiêu thụ hết.
\begin{enumerate}
    \item Gọi $x$ và $y$ lần lượt là số đơn vị sản phẩm loại $A$ và loại $B$ mà xưởng sản xuất trong một ngày. Hãy thiết lập hệ bất phương trình bậc nhất hai ẩn mô tả các điều kiện ràng buộc đối với $x$ và $y$.
    \item Xác định số đơn vị sản phẩm mỗi loại cần sản xuất để xưởng đạt lợi nhuận lớn nhất và tính số tiền lợi nhuận lớn nhất đó.
\end{enumerate}
\end{enumerate}
\newpage

% =========================================================================
% TRANG 4 TRỞ ĐI: BẢNG ĐÁP ÁN TỔNG HỢP VÀ LỜI GIẢI CHI TIẾT
% =========================================================================
\parttitle{BẢNG ĐÁP ÁN TỔNG HỢP}

\vspace{0.1cm}
\textbf{1. Phần I: Trắc nghiệm nhiều lựa chọn}
\begin{center}
{\setlength{\tabcolsep}{2.5pt}\small
\begin{tabular}{|c|c|c|c|c|c|c|c|c|c|c|c|}
\hline
\textbf{Câu 1} & \textbf{Câu 2} & \textbf{Câu 3} & \textbf{Câu 4} & \textbf{Câu 5} & \textbf{Câu 6} & \textbf{Câu 7} & \textbf{Câu 8} & \textbf{Câu 9} & \textbf{Câu 10} & \textbf{Câu 11} & \textbf{Câu 12} \\
\hline
\textbf{A} & \textbf{B} & \textbf{C} & \textbf{D} & \textbf{A} & \textbf{B} & \textbf{C} & \textbf{D} & \textbf{A} & \textbf{B} & \textbf{C} & \textbf{D} \\
\hline
\end{tabular}}
\end{center}

\vspace{0.2cm}
\textbf{2. Phần II: Trắc nghiệm đúng/sai}
\begin{center}
{\small
\begin{tabular}{|c|c|c|c|c|}
\hline
\textbf{Câu} & \textbf{Lệnh hỏi a)} & \textbf{Lệnh hỏi b)} & \textbf{Lệnh hỏi c)} & \textbf{Lệnh hỏi d)} \\
\hline
\textbf{1} & \textbf{Đ} & \textbf{S} & \textbf{Đ} & \textbf{Đ} \\
\hline
\textbf{2} & \textbf{Đ} & \textbf{Đ} & \textbf{Đ} & \textbf{S} \\
\hline
\end{tabular}}
\end{center}

\vspace{0.2cm}
\textbf{3. Phần III: Trắc nghiệm trả lời ngắn}
\begin{center}
{\small
\begin{tabular}{|c|c|c|c|c|}
\hline
\textbf{Câu} & \textbf{Câu 1} & \textbf{Câu 2} & \textbf{Câu 3} & \textbf{Câu 4} \\
\hline
\textbf{Đáp số} & $6$ & $82$ & $12$ & $7$ \\
\hline
\end{tabular}}
\end{center}

\vspace{0.3cm}
\parttitle{LỜI GIẢI CHI TIẾT}

\textbf{\large PHẦN I. Trắc nghiệm nhiều lựa chọn}
\begin{enumerate}[label=\textbf{Câu \arabic*.},leftmargin=1.1cm]
\item \textbf{Đáp án A.}
Câu ``$17$ là số nguyên tố'' là một câu khẳng định có tính đúng sai xác định (và là khẳng định đúng), do đó nó là một mệnh đề toán học. Các câu ở phương án B, C, D lần lượt là câu cảm thán, câu hỏi và câu cầu khiến nên không phải mệnh đề.

\item \textbf{Đáp án B.}
Phủ định của mệnh đề chứa kí hiệu với mọi $\forall x \in \mathbb{R},\, P(x)$ là $\exists x \in \mathbb{R},\, \overline{P(x)}$. Do đó, phủ định của $\forall x \in \mathbb{R},\, x^2 - 3x + 2 > 0$ là $\exists x \in \mathbb{R},\, x^2 - 3x + 2 \le 0$.

\item \textbf{Đáp án C.}
Phương trình $(2x - 1)(x^2 - 4) = 0 \iff \left[\begin{array}{l} 2x - 1 = 0 \\ x^2 - 4 = 0 \end{array}\right. \iff \left[\begin{array}{l} x = \dfrac{1}{2} \\ x = 2 \\ x = -2 \end{array}\right.$.
Do $x \in \mathbb{N}$ nên ta chỉ nhận nghiệm tự nhiên duy nhất là $x = 2$. Vậy $A = \{2\}$.

\item \textbf{Đáp án D.}
Một tập hợp hữu hạn có $n$ phần tử thì có tổng cộng $2^n$ tập hợp con. Tập hợp $X = \{a;\, b;\, c\}$ có $3$ phần tử nên số tập con của $X$ là $2^3 = 8$.

\item \textbf{Đáp án A.}
Bất phương trình $2x - 3y + 5 > 0$ có dạng bậc nhất đối với hai ẩn $x$ và $y$. Các bất phương trình ở phương án B, C có chứa $x^2, y^2$, phương án D chứa tích $xy$ nên không phải là bất phương trình bậc nhất hai ẩn.

\item \textbf{Đáp án B.}
Thay tọa độ điểm $(2;\, 1)$ vào vế trái của bất phương trình:
$2(2) - 1 + 3 = 4 - 1 + 3 = 6 \ge 0$ (thỏa mãn).
Các điểm $(1;\, 6)$, $(0;\, 4)$, $(-3;\, 1)$ khi thay vào đều cho giá trị âm.

\item \textbf{Đáp án C.}
Thay cặp số $(0;\, 1)$ vào hệ:
$\begin{cases} 0 + 1 - 2 = -1 \le 0 \\ 2(0) - 3(1) + 6 = 3 > 0 \end{cases}$ (thỏa mãn cả hai bất phương trình).
Do đó $(0;\, 1)$ là nghiệm của hệ bất phương trình.

\item \textbf{Đáp án D.}
Bờ của miền nghiệm của bất phương trình $3x - 2y + 6 \ge 0$ chính là đường thẳng có phương trình $3x - 2y + 6 = 0$.

\item \textbf{Đáp án A.}
Ta có $\cos 120^\circ = \cos(180^\circ - 60^\circ) = -\cos 60^\circ = -\dfrac{1}{2}$.

\item \textbf{Đáp án B.}
Áp dụng định lý côsin trong tam giác $ABC$:
\[
BC^2 = AB^2 + AC^2 - 2AB \cdot AC \cdot \cos A = 5^2 + 8^2 - 2 \cdot 5 \cdot 8 \cdot \cos 60^\circ = 25 + 64 - 80 \cdot \dfrac{1}{2} = 49.
\]
Suy ra $BC = \sqrt{49} = 7$.

\item \textbf{Đáp án C.}
Ta nhận thấy $a^2 + b^2 = 6^2 + 8^2 = 36 + 64 = 100 = 10^2 = c^2$.
Theo định lý Pythagore đảo, tam giác $ABC$ vuông tại $C$ với cạnh huyền $c = 10$.
Bán kính đường tròn ngoại tiếp tam giác vuông bằng nửa cạnh huyền: $R = \dfrac{c}{2} = \dfrac{10}{2} = 5$.

\item \textbf{Đáp án D.}
Theo công thức tính độ dài đường trung tuyến trong tam giác, ta có $m_a^2 = \dfrac{2(b^2 + c^2) - a^2}{4}$.
\end{enumerate}

\vspace{0.3cm}
\textbf{\large PHẦN II. Trắc nghiệm đúng/sai}
\begin{enumerate}[label=\textbf{Câu \arabic*.},leftmargin=1.1cm]
\item 
Cho hai tập hợp $A = [-2;\, 4]$ và $B = (1;\, +\infty)$.
\begin{enumerate}
    \item \textbf{Đúng.} Vì $-2 \le 0 \le 4$ nên $0 \in [-2;\, 4] = A$.
    \item \textbf{Sai.} Tập hợp $B = \{x \in \mathbb{R} \mid x > 1\}$ là khoảng $(1;\, +\infty)$, không phải nửa khoảng $[1;\, +\infty)$ (vì không chứa điểm $1$).
    \item \textbf{Đúng.} Giao của hai tập hợp $A \cap B = [-2;\, 4] \cap (1;\, +\infty) = (1;\, 4]$.
    \item \textbf{Đúng.} Hiệu của hai tập hợp $A \setminus B$ là tập hợp gồm các phần tử thuộc $A$ nhưng không thuộc $B$. Vì $B = (1;\, +\infty)$ nên $A \setminus B = [-2;\, 4] \setminus (1;\, +\infty) = [-2;\, 1]$.
\end{enumerate}

\item 
Cho tam giác $ABC$ có độ dài ba cạnh $a = 7$, $b = 5$, $c = 6$.
\begin{enumerate}
    \item \textbf{Đúng.} Nửa chu vi tam giác là $p = \dfrac{a + b + c}{2} = \dfrac{7 + 5 + 6}{2} = 9$.
    \item \textbf{Đúng.} Áp dụng hệ quả định lý côsin: $\cos A = \dfrac{b^2 + c^2 - a^2}{2bc} = \dfrac{5^2 + 6^2 - 7^2}{2 \cdot 5 \cdot 6} = \dfrac{12}{60} = \dfrac{1}{5} > 0$. Vì $\cos A > 0$ nên góc $A$ là góc nhọn.
    \item \textbf{Đúng.} Theo tính toán ở trên, $\cos A = \dfrac{1}{5}$.
    \item \textbf{Sai.} Áp dụng công thức Heron, diện tích tam giác $ABC$ là:
    \[
    S = \sqrt{p(p - a)(p - b)(p - c)} = \sqrt{9(9 - 7)(9 - 5)(9 - 6)} = \sqrt{9 \cdot 2 \cdot 4 \cdot 3} = \sqrt{216} = 6\sqrt{6}.
    \]
    Do đó mệnh đề khẳng định diện tích bằng $12\sqrt{6}$ là sai.
\end{enumerate}
\end{enumerate}

\vspace{0.3cm}
\textbf{\large PHẦN III. Trắc nghiệm trả lời ngắn}
\begin{enumerate}[label=\textbf{Câu \arabic*.},leftmargin=1.1cm]
\item \textbf{Đáp số: $6$.}
Đường thẳng $d$ đi qua hai điểm $(4;\, 0)$ và $(0;\, 2)$ nên có phương trình đoạn chắn $\dfrac{x}{4} + \dfrac{y}{2} = 1 \iff x + 2y = 4$.
Điểm gốc tọa độ $O(0;\, 0)$ thuộc miền không bị gạch chéo. Thay $(0;\, 0)$ vào biểu thức $x + 2y$ ta được $0 + 2(0) = 0 \le 4$.
Do đó, miền nghiệm không bị gạch chéo (kể cả bờ $d$) là miền nghiệm của bất phương trình $x + 2y \le 4$.
Đồng nhất với dạng $x + by \le c$, ta có $b = 2$ và $c = 4$.
Vậy $T = b + c = 2 + 4 = 6$.

\item \textbf{Đáp số: $82$.}
Xét tam giác $ABC$, ta có tổng ba góc trong tam giác bằng $180^\circ$:
\[
\widehat{B} = 180^\circ - (\widehat{A} + \widehat{C}) = 180^\circ - (60^\circ + 75^\circ) = 45^\circ.
\]
Áp dụng định lý sin trong tam giác $ABC$:
\[
\dfrac{AB}{\sin C} = \dfrac{AC}{\sin B} \implies AB = \dfrac{AC \cdot \sin C}{\sin B} = \dfrac{60 \cdot \sin 75^\circ}{\sin 45^\circ}.
\]
Ta có $\sin 75^\circ = \sin(45^\circ + 30^\circ) = \dfrac{\sqrt{6} + \sqrt{2}}{4}$ và $\sin 45^\circ = \dfrac{\sqrt{2}}{2}$.
Suy ra:
\[
AB = \dfrac{60 \cdot \dfrac{\sqrt{6} + \sqrt{2}}{4}}{\dfrac{\sqrt{2}}{2}} = 30(\sqrt{3} + 1) \approx 30 \times 2{,}73205 = 81{,}9615\text{ m}.
\]
Làm tròn đến hàng đơn vị, khoảng cách giữa $A$ và $B$ là $82\text{ m}$.

\item \textbf{Đáp số: $12$.}
Gọi $A$ là tập hợp học sinh thích môn Bóng đá, $B$ là tập hợp học sinh thích môn Bóng rổ.
Theo giả thiết: $|A| = 28$, $|B| = 23$.
Số học sinh thích ít nhất một môn thể thao là:
\[
|A \cup B| = 45 - 6 = 39.
\]
Áp dụng công thức cộng số phần tử của hai tập hợp:
\[
|A \cup B| = |A| + |B| - |A \cap B| \implies |A \cap B| = |A| + |B| - |A \cup B| = 28 + 23 - 39 = 12.
\]
Vậy có $12$ học sinh thích cả hai môn Bóng đá và Bóng rổ.

\item \textbf{Đáp số: $7$.}
Áp dụng định lý côsin trong tam giác $ABC$:
\[
BC^2 = AB^2 + AC^2 - 2AB \cdot AC \cdot \cos A = 3^2 + 5^2 - 2 \cdot 3 \cdot 5 \cdot \cos 120^\circ.
\]
Vì $\cos 120^\circ = -\dfrac{1}{2}$ nên:
\[
BC^2 = 9 + 25 - 30 \cdot \left(-\dfrac{1}{2}\right) = 34 + 15 = 49 \implies BC = \sqrt{49} = 7.
\]
\end{enumerate}

\vspace{0.3cm}
\textbf{\large PHẦN IV. Tự luận}
\begin{enumerate}[label=\textbf{Câu \arabic*.},leftmargin=1.1cm]
\item 
Vì $90^\circ < \alpha < 180^\circ$ nên $\sin \alpha > 0$.
Áp dụng công thức lượng giác cơ bản $\sin^2 \alpha + \cos^2 \alpha = 1$, ta có:
\[
\sin \alpha = \sqrt{1 - \cos^2 \alpha} = \sqrt{1 - \left(-\dfrac{3}{5}\right)^2} = \sqrt{1 - \dfrac{9}{25}} = \sqrt{\dfrac{16}{25}} = \dfrac{4}{5}.
\]
Thay $\sin \alpha = \dfrac{4}{5}$ và $\cos \alpha = -\dfrac{3}{5}$ vào biểu thức $P$:
\[
P = \dfrac{2\sin \alpha - 3\cos \alpha}{\sin \alpha + \cos \alpha} = \dfrac{2 \cdot \left(\dfrac{4}{5}\right) - 3 \cdot \left(-\dfrac{3}{5}\right)}{\dfrac{4}{5} + \left(-\dfrac{3}{5}\right)} = \dfrac{\dfrac{8}{5} + \dfrac{9}{5}}{\dfrac{1}{5}} = \dfrac{\dfrac{17}{5}}{\dfrac{1}{5}} = 17.
\]
Vậy giá trị của biểu thức $P$ bằng $17$.

\item 
Cho hai tập hợp $A = [-3;\, 2)$ và $B = (0;\, 5]$.
\begin{enumerate}
    \item Xác định các tập hợp:
    \begin{itemize}
        \item $A \cup B = [-3;\, 2) \cup (0;\, 5] = [-3;\, 5]$.
        \item $A \cap B = [-3;\, 2) \cap (0;\, 5] = (0;\, 2)$.
        \item $A \setminus B = [-3;\, 2) \setminus (0;\, 5] = [-3;\, 0]$.
    \end{itemize}
    \item Biểu diễn tập hợp $A \cap B = (0;\, 2)$ trên trục số:
    \begin{center}
    \begin{tikzpicture}[scale=1.0, >=Stealth]
        % Trục số
        \draw[line width=1.1pt, ->] (-2,0) -- (4.5,0) node[right, font=\small\bfseries] {$x$};
        
        % Gạch bỏ phần ngoài (0; 2)
        \foreach \x in {-1.8,-1.5,...,-0.1}
            \draw[line width=0.7pt, gray] (\x-0.1, -0.2) -- (\x+0.1, 0.2);
        \foreach \x in {2.1,2.4,...,4.2}
            \draw[line width=0.7pt, gray] (\x-0.1, -0.2) -- (\x+0.1, 0.2);
            
        % Điểm mút ngoặc tròn
        \node[font=\large\bfseries, color=ink] at (0, 0) {$($};
        \node[font=\large\bfseries, color=ink] at (2, 0) {$)$};
        
        % Nhãn tọa độ
        \node[below=3pt, font=\small\bfseries] at (0,0) {$0$};
        \node[below=3pt, font=\small\bfseries] at (2,0) {$2$};
    \end{tikzpicture}
    \end{center}
\end{enumerate}

\item 
Vì độ dài của đoạn $A$ là $(m + 3) - (m - 1) = 4 > 0$ nên tập hợp $A = [m - 1;\, m + 3]$ luôn khác rỗng với mọi $m \in \mathbb{R}$.
Để $A$ là tập con của $B = (-2;\, 4)$, tức là mọi phần tử của $A$ đều thuộc khoảng $(-2;\, 4)$, điều kiện cần và đủ là:
\[
\begin{cases}
m - 1 > -2 \\
m + 3 < 4
\end{cases}
\iff
\begin{cases}
m > -1 \\
m < 1
\end{cases}
\iff -1 < m < 1.
\]
Vậy tất cả các giá trị thực của tham số $m$ thỏa mãn yêu cầu bài toán là $m \in (-1;\, 1)$.

\item 
\begin{enumerate}
    \item Gọi $x \ge 0$ và $y \ge 0$ lần lượt là số đơn vị sản phẩm loại $A$ và loại $B$ mà xưởng sản xuất trong một ngày.
    \begin{itemize}
        \item Khối lượng thép cần dùng là: $2x + y$ (kg). Vì lượng thép sẵn có tối đa là $10\text{ kg}$ nên ta có bất phương trình:
        \[
        2x + y \le 10.
        \]
        \item Thời gian máy hoạt động là: $x + 2y$ (giờ). Vì máy hoạt động tối đa $8\text{ giờ}$ mỗi ngày nên ta có bất phương trình:
        \[
        x + 2y \le 8.
        \]
    \end{itemize}
    Vậy hệ bất phương trình bậc nhất hai ẩn mô tả các điều kiện ràng buộc đối với $x$ và $y$ là:
    \[
    \begin{cases}
    2x + y \le 10 \\
    x + 2y \le 8 \\
    x \ge 0 \\
    y \ge 0.
    \end{cases}
    \]
    \item Tổng lợi nhuận xưởng thu được mỗi ngày là biểu thức mục tiêu:
    \[
    F(x;\, y) = 400x + 300y\quad \text{(nghìn đồng)}.
    \]
    Miền nghiệm của hệ bất phương trình ràng buộc là miền tứ giác $OABC$ (kể cả biên) trong mặt phẳng tọa độ $Oxy$ dưới đây:
    
    \begin{center}
    \begin{tikzpicture}[scale=0.88, >=Stealth]
        % Trục tọa độ chống đè chữ
        \draw[line width=1.1pt, ->] (-0.8,0) -- (6.5,0) node[right, font=\small\bfseries] {$x$};
        \draw[line width=1.1pt, ->] (0,-0.8) -- (0,5.5) node[above, font=\small\bfseries] {$y$};
        \node[below left=2pt, font=\small\bfseries] at (0,0) {$O$};
        
        % Miền tứ giác OABC
        \fill[fill=ink!15] (0,0) -- (5,0) -- (4,2) -- (0,4) -- cycle;
        
        % Đường thẳng d1: 2x + y = 10 đi qua (5, 0) và (4, 2)
        \draw[line width=1.2pt, ink] (2.2, 5.6) -- (5.6, -1.2);
        \node[above right=1pt, font=\small\bfseries] at (2.2, 5.6) {$d_1\colon 2x + y = 10$};
        
        % Đường thẳng d2: x + 2y = 8 đi qua (0, 4) và (4, 2)
        \draw[line width=1.2pt, blue!70!black] (-0.8, 4.4) -- (6.4, 0.8);
        \node[above right=1pt, font=\small\bfseries] at (6.0, 1.0) {$d_2\colon x + 2y = 8$};
        
        % Đường gióng tọa độ điểm B(4; 2)
        \draw[dashed, line width=0.8pt, gray] (4,0) -- (4,2) -- (0,2);
        
        % Tọa độ các đỉnh
        \filldraw[black] (0,0) circle (2.2pt);
        \filldraw[black] (5,0) circle (2.4pt) node[below=2pt, font=\small\bfseries] {$5$} node[above right=2pt, font=\small\bfseries] {$A$};
        \filldraw[black] (4,2) circle (2.4pt) node[above right=2pt, font=\small\bfseries] {$B(4;\, 2)$};
        \filldraw[black] (0,4) circle (2.4pt) node[left=2pt, font=\small\bfseries] {$4$} node[above right=2pt, font=\small\bfseries] {$C$};
        
        % Các vạch số trên trục
        \filldraw[black] (4,0) circle (1.8pt) node[below=2pt, font=\small\bfseries] {$4$};
        \filldraw[black] (0,2) circle (1.8pt) node[left=2pt, font=\small\bfseries] {$2$};
        
        % Nhãn miền nghiệm
        \node[font=\small\bfseries\color{ink}] at (1.8, 1.4) {Miền nghiệm};
    \end{tikzpicture}
    \end{center}
    
    Tọa độ các đỉnh của tứ giác $OABC$ là:
    \begin{itemize}
        \item $O(0;\, 0)$, $A(5;\, 0)$, $C(0;\, 4)$.
        \item $B(4;\, 2)$ là giao điểm của hai đường thẳng $d_1$ và $d_2$ (nghiệm của hệ $\begin{cases} 2x + y = 10 \\ x + 2y = 8 \end{cases}$).
    \end{itemize}
    
    Tính giá trị của biểu thức lợi nhuận $F(x;\, y)$ tại các đỉnh:
    \begin{itemize}
        \item Tại $O(0;\, 0)$: $F(0;\, 0) = 400(0) + 300(0) = 0$ (nghìn đồng).
        \item Tại $A(5;\, 0)$: $F(5;\, 0) = 400(5) + 300(0) = 2\,000$ (nghìn đồng).
        \item Tại $B(4;\, 2)$: $F(4;\, 2) = 400(4) + 300(2) = 1\,600 + 600 = 2\,200$ (nghìn đồng).
        \item Tại $C(0;\, 4)$: $F(0;\, 4) = 400(0) + 300(4) = 1\,200$ (nghìn đồng).
    \end{itemize}
    So sánh giá trị tại các đỉnh, ta thấy giá trị lớn nhất của $F(x;\, y)$ là $2\,200$ nghìn đồng, đạt được khi $x = 4$ và $y = 2$.
    
    \textbf{Kết luận:} Xưởng cần sản xuất $4$ đơn vị sản phẩm loại $A$ và $2$ đơn vị sản phẩm loại $B$ mỗi ngày để đạt lợi nhuận lớn nhất là $2\,200\,000$ đồng (hai triệu hai trăm nghìn đồng).
\end{enumerate}
\end{enumerate}

\end{document}
